% ============================================================================
% CAPÍTULO 3: MODELOS DE REDES ESTÁTICAS
% Bioinformática para Biologia de Sistemas
% ============================================================================

% ============================================================================
% TEMPLATE BASE PARA APRESENTAÇÕES EM BEAMER
% Bioinformática para Biologia de Sistemas
% ============================================================================

\documentclass[12pt, aspectratio=169]{beamer}

% ============================================================================
% PACOTES ESSENCIAIS
% ============================================================================
\usepackage[utf8]{inputenc}
\usepackage[T1]{fontenc}
\usepackage[brazilian]{babel}
\usepackage{amsmath, amssymb, amsthm}
\usepackage{graphicx}
\usepackage{xcolor}
\usepackage{tikz}
\usetikzlibrary{arrows.meta, shapes, positioning, calc, decorations.pathreplacing, decorations.shapes, decorations.pathmorphing}
\usepackage{listings}
\usepackage{booktabs}
\usepackage{multirow}
\usepackage{hyperref}
\usepackage{chemfig}
\usepackage{chemarr}

% ============================================================================
% CONFIGURAÇÃO DO TEMA
% ============================================================================
\usetheme{Madrid}
\usecolortheme{default}

% Cores personalizadas
\definecolor{azulescuro}{RGB}{0, 51, 102}
\definecolor{azulclaro}{RGB}{51, 153, 255}
\definecolor{verdeacido}{RGB}{102, 204, 0}
\definecolor{laranjacel}{RGB}{255, 153, 51}
\definecolor{roxodna}{RGB}{153, 51, 255}

\setbeamercolor{structure}{fg=azulescuro}
\setbeamercolor{palette primary}{bg=azulescuro,fg=white}
\setbeamercolor{palette secondary}{bg=azulclaro,fg=white}
\setbeamercolor{palette tertiary}{bg=azulescuro,fg=white}
\setbeamercolor{palette quaternary}{bg=azulclaro,fg=white}
\setbeamercolor{block title}{bg=azulescuro,fg=white}
\setbeamercolor{block body}{bg=azulclaro!10,fg=black}
\setbeamercolor{block title alerted}{bg=red!80,fg=white}
\setbeamercolor{block body alerted}{bg=red!10,fg=black}
\setbeamercolor{block title example}{bg=verdeacido!80,fg=white}
\setbeamercolor{block body example}{bg=verdeacido!10,fg=black}

% ============================================================================
% CONFIGURAÇÃO DE CÓDIGO
% ============================================================================
\lstset{
    language=Python,
    basicstyle=\ttfamily\small,
    keywordstyle=\color{azulescuro}\bfseries,
    commentstyle=\color{verdeacido},
    stringstyle=\color{laranjacel},
    numbers=left,
    numberstyle=\tiny\color{gray},
    stepnumber=1,
    numbersep=5pt,
    backgroundcolor=\color{white},
    showspaces=false,
    showstringspaces=false,
    showtabs=false,
    frame=single,
    rulecolor=\color{black},
    tabsize=2,
    captionpos=b,
    breaklines=true,
    breakatwhitespace=false,
    escapeinside={\%*}{*)},
    xleftmargin=10pt,
    xrightmargin=10pt
}

% Definir estilo para R
\lstdefinestyle{Rstyle}{
    language=R,
    basicstyle=\ttfamily\small,
    keywordstyle=\color{azulescuro}\bfseries,
    commentstyle=\color{verdeacido},
    stringstyle=\color{laranjacel}
}

% Definir estilo para MATLAB
\lstdefinestyle{MATLABstyle}{
    language=Matlab,
    basicstyle=\ttfamily\small,
    keywordstyle=\color{azulescuro}\bfseries,
    commentstyle=\color{verdeacido},
    stringstyle=\color{laranjacel}
}

% ============================================================================
% COMANDOS PERSONALIZADOS
% ============================================================================

% Comando para equações destacadas
\newcommand{\destaque}[1]{%
    \begin{alertblock}{}
        \begin{center}
            #1
        \end{center}
    \end{alertblock}
}

% Comando para definições
\newcommand{\definicao}[2]{%
    \begin{block}{Definição: #1}
        #2
    \end{block}
}

% Comando para exemplos
\newcommand{\exemplo}[2]{%
    \begin{exampleblock}{Exemplo: #1}
        #2
    \end{exampleblock}
}

% Comando para observações importantes
\newcommand{\importante}[1]{%
    \begin{alertblock}{Importante!}
        #1
    \end{alertblock}
}

% Comandos matemáticos úteis
\newcommand{\R}{\mathbb{R}}
\newcommand{\N}{\mathbb{N}}
\newcommand{\Z}{\mathbb{Z}}
\newcommand{\dif}{\mathrm{d}}
\newcommand{\parcial}[2]{\frac{\partial #1}{\partial #2}}
\newcommand{\derivada}[2]{\frac{\dif #1}{\dif #2}}
\newcommand{\vect}[1]{\boldsymbol{#1}}

% Comandos biológicos
\newcommand{\gene}[1]{\textit{#1}}
\newcommand{\proteina}[1]{\textsc{#1}}
\newcommand{\especie}[1]{\textit{#1}}

% ============================================================================
% CONFIGURAÇÃO DE RODAPÉ
% ============================================================================
\setbeamertemplate{footline}{
  \leavevmode%
  \hbox{%
  \begin{beamercolorbox}[wd=.33\paperwidth,ht=2.25ex,dp=1ex,center]{author in head/foot}%
    \usebeamerfont{author in head/foot}\insertshortauthor
  \end{beamercolorbox}%
  \begin{beamercolorbox}[wd=.34\paperwidth,ht=2.25ex,dp=1ex,center]{title in head/foot}%
    \usebeamerfont{title in head/foot}\insertshorttitle
  \end{beamercolorbox}%
  \begin{beamercolorbox}[wd=.33\paperwidth,ht=2.25ex,dp=1ex,right]{date in head/foot}%
    \usebeamerfont{date in head/foot}\insertshortdate{}\hspace*{2em}
    \insertframenumber{} / \inserttotalframenumber\hspace*{2ex}
  \end{beamercolorbox}}%
  \vskip0pt%
}

% ============================================================================
% CONFIGURAÇÃO DE NAVEGAÇÃO
% ============================================================================
\setbeamertemplate{navigation symbols}{}

% ============================================================================
% CONFIGURAÇÃO DE ITEMIZE/ENUMERATE
% ============================================================================
\setbeamertemplate{itemize items}[circle]
\setbeamertemplate{enumerate items}[default]

% ============================================================================
% FIM DO TEMPLATE
% ============================================================================


% ============================================================================
% INFORMAÇÕES DO DOCUMENTO
% ============================================================================
\title[Cap. 3: Modelos de Redes Estáticas]{Capítulo 3: Modelos de Redes Estáticas}
\subtitle{Bioinformática para Biologia de Sistemas}
\author{Ney Lemke}
\institute{DBF-Unesp}
\date{\today}

\begin{document}

% ============================================================================
% SLIDE DE TÍTULO
% ============================================================================
\begin{frame}
\titlepage
\end{frame}

% ============================================================================
% SUMÁRIO
% ============================================================================
\begin{frame}{Sumário}
\tableofcontents
\end{frame}

% ============================================================================
% SEÇÃO 1: INTRODUÇÃO
% ============================================================================
\section{Introdução}

\begin{frame}[shrink=10]{Visão Geral do Capítulo (1/2)}
\begin{block}{Objetivos de Aprendizagem}
\begin{itemize}
    \item Compreender os fundamentos de teoria dos grafos
    \item Analisar propriedades topológicas de redes biológicas
    \item Explorar modelos de redes small-world e scale-free
\end{itemize}
\end{block}
\end{frame}

\begin{frame}[shrink=10]{Visão Geral do Capítulo (2/2)}
\begin{exampleblock}{Por que Estudar Redes Biológicas?}
\begin{itemize}
    \item Entender organização e função de sistemas biológicos
    \item Identificar alvos terapêuticos e biomarcadores
    \item Predizer efeitos de perturbações no sistema
    \item Descobrir princípios de design evolutivo
\end{itemize}
\end{exampleblock}
\end{frame}

\begin{frame}[shrink=10]{O que são Redes Biológicas?}
\definicao{Rede Biológica}{
Uma rede biológica é uma representação matemática de interações entre componentes biológicos (genes, proteínas, metabólitos) como um grafo, onde nós representam entidades e arestas representam relações.
}

\vspace{0.3cm}

\begin{columns}
\column{0.5\textwidth}
\textbf{Características:}
\begin{itemize}
    \item Estrutura modular
    \item Hierarquia organizacional
    \item Robustez a perturbações
    \item Propriedades emergentes
\end{itemize}

\column{0.5\textwidth}
\begin{center}
\begin{tikzpicture}[scale=0.6]
    \foreach \i in {1,...,6} {
        \node[circle, draw, fill=azulclaro!30] (n\i) at ({60*\i}:2) {\i};
    }
    \draw[thick, azulescuro] (n1) -- (n2);
    \draw[thick, azulescuro] (n2) -- (n3);
    \draw[thick, azulescuro] (n3) -- (n4);
    \draw[thick, azulescuro] (n4) -- (n5);
    \draw[thick, azulescuro] (n5) -- (n6);
    \draw[thick, azulescuro] (n6) -- (n1);
    \draw[thick, verdeacido] (n1) -- (n4);
    \draw[thick, verdeacido] (n2) -- (n5);
    \draw[thick, verdeacido] (n3) -- (n6);
\end{tikzpicture}
\end{center}
\end{columns}
\end{frame}

\begin{frame}[shrink=5]{Tipos de Redes Biológicas (1/2)}
\begin{columns}
\column{0.5\textwidth}
\begin{block}{Redes de Interação Proteica (PPI)}
\begin{itemize}
    \item Nós: proteínas
    \item Arestas: interações físicas
    \item Não-direcionadas
    \item Ex: complexos proteicos
\end{itemize}
\end{block}

\column{0.5\textwidth}
\begin{block}{Redes Metabólicas}
\begin{itemize}
    \item Nós: metabólitos ou reações
    \item Arestas: conversões enzimáticas
    \item Direcionadas
    \item Ex: glicólise, ciclo de Krebs
\end{itemize}
\end{block}
\end{columns}
\end{frame}

\begin{frame}[shrink=5]{Tipos de Redes Biológicas (2/2)}
\begin{columns}
\column{0.5\textwidth}
\begin{block}{Redes Regulatórias Gênicas}
\begin{itemize}
    \item Nós: genes/fatores de transcrição
    \item Arestas: regulação transcricional
    \item Direcionadas (ativação/repressão)
    \item Ex: circuitos genéticos
\end{itemize}
\end{block}

\column{0.5\textwidth}
\begin{block}{Redes de Sinalização}
\begin{itemize}
    \item Nós: proteínas sinalizadoras
    \item Arestas: cascatas de sinalização
    \item Direcionadas
    \item Ex: vias MAPK, PI3K-AKT
\end{itemize}
\end{block}
\end{columns}
\end{frame}

% ============================================================================
% SEÇÃO 2: FUNDAMENTOS DE TEORIA DOS GRAFOS
% ============================================================================
\section{Fundamentos de Teoria dos Grafos}

\begin{frame}[shrink=5]{Conceitos Básicos de Grafos}
\definicao{Grafo}{
Um grafo $G = (V, E)$ é composto por:
\begin{itemize}
    \item $V$: conjunto de vértices (nós)
    \item $E$: conjunto de arestas (conexões) entre vértices
\end{itemize}
}

\vspace{0.3cm}

\begin{columns}
\column{0.5\textwidth}
\begin{exampleblock}{Grafo Não-Direcionado}
Arestas sem direção: $(u, v) = (v, u)$
\begin{center}
\begin{tikzpicture}[scale=0.7]
    \node[circle, draw, fill=azulclaro!30] (A) at (0,0) {A};
    \node[circle, draw, fill=azulclaro!30] (B) at (2,0) {B};
    \node[circle, draw, fill=azulclaro!30] (C) at (1,1.5) {C};
    \draw[thick] (A) -- (B);
    \draw[thick] (B) -- (C);
    \draw[thick] (C) -- (A);
\end{tikzpicture}
\end{center}
\end{exampleblock}

\column{0.5\textwidth}
\begin{exampleblock}{Grafo Direcionado (Dígrafo)}
Arestas com direção: $(u, v) \neq (v, u)$
\begin{center}
\begin{tikzpicture}[scale=0.7]
    \node[circle, draw, fill=verdeacido!30] (A) at (0,0) {A};
    \node[circle, draw, fill=verdeacido!30] (B) at (2,0) {B};
    \node[circle, draw, fill=verdeacido!30] (C) at (1,1.5) {C};
    \draw[->, thick] (A) -- (B);
    \draw[->, thick] (B) -- (C);
    \draw[->, thick] (C) -- (A);
\end{tikzpicture}
\end{center}
\end{exampleblock}
\end{columns}
\end{frame}

\begin{frame}[shrink=15]{Representação de Grafos: Matriz de Adjacência}
\definicao{Matriz de Adjacência}{
Matriz $A$ de dimensão $n \times n$ onde $n = |V|$: $A_{ij} = 1$ se existe aresta de $i$ para $j$, e $A_{ij} = 0$ caso contrário.
}

\begin{columns}
\column{0.45\textwidth}
\begin{center}
\begin{tikzpicture}[scale=0.9]
    \node[circle, draw, fill=azulclaro!30] (1) at (0,0) {1};
    \node[circle, draw, fill=azulclaro!30] (2) at (2.5,0) {2};
    \node[circle, draw, fill=azulclaro!30] (3) at (0,2.5) {3};
    \node[circle, draw, fill=azulclaro!30] (4) at (2.5,2.5) {4};
    \draw[thick] (1) -- (2);
    \draw[thick] (1) -- (3);
    \draw[thick] (2) -- (4);
    \draw[thick] (3) -- (4);
\end{tikzpicture}
\end{center}

\column{0.55\textwidth}
\begin{center}
\textbf{Matriz de Adjacência:}

$A = \begin{bmatrix}
0 & 1 & 1 & 0 \\
1 & 0 & 0 & 1 \\
1 & 0 & 0 & 1 \\
0 & 1 & 1 & 0
\end{bmatrix}$
\end{center}

\begin{exampleblock}{Propriedades}
\begin{itemize}
    \item Simétrica para grafos não-direcionados
    \item Diagonal zero (sem auto-loops)
    \item $\sum_j A_{ij} = k_i$ \ (grau do vértice $i$)
\end{itemize}
\end{exampleblock}
\end{columns}
\end{frame}

\begin{frame}[shrink=5]{Representação de Grafos: Matriz de Incidência}
\definicao{Matriz de Incidência}{
Matriz $B$ de dimensão $n \times m$ onde $n = |V|$ e $m = |E|$:
\begin{equation*}
B_{ie} = \begin{cases}
1 & \text{se vértice } i \text{ é origem da aresta } e \\
-1 & \text{se vértice } i \text{ é destino da aresta } e \\
0 & \text{caso contrário}
\end{cases}
\end{equation*}
}

\vspace{0.3cm}

\begin{exampleblock}{Vantagens}
\begin{itemize}
    \item Útil para representar grafos esparsos
    \item Facilita análise de fluxos em redes metabólicas
    \item Cada coluna representa uma reação bioquímica
\end{itemize}
\end{exampleblock}
\end{frame}

\begin{frame}[shrink=20]{Grau de um Vértice}
\definicao{Grau (Degree)}{
O grau de um vértice é o número de arestas conectadas a ele.
\begin{itemize}
    \item \textbf{Grafo não-direcionado:} $k_i = \sum_{j} A_{ij}$
    \item \textbf{Grafo direcionado:}
    \begin{itemize}
        \item In-degree: $k_i^{in} = \sum_{j} A_{ji}$ (arestas que chegam)
        \item Out-degree: $k_i^{out} = \sum_{j} A_{ij}$ (arestas que saem)
    \end{itemize}
\end{itemize}
}

\begin{columns}
\column{0.5\textwidth}
\begin{center}
\begin{tikzpicture}[scale=0.8]
    \node[circle, draw, fill=laranjacel!50] (A) at (1,1) {A};
    \node[circle, draw, fill=azulclaro!30] (B) at (0,0) {B};
    \node[circle, draw, fill=azulclaro!30] (C) at (2,0) {C};
    \node[circle, draw, fill=azulclaro!30] (D) at (0,2) {D};
    \node[circle, draw, fill=azulclaro!30] (E) at (2,2) {E};
    \draw[thick] (A) -- (B);
    \draw[thick] (A) -- (C);
    \draw[thick] (A) -- (D);
    \draw[thick] (A) -- (E);
\end{tikzpicture}
\end{center}

\column{0.5\textwidth}
\textbf{Nó A:} $k_A = 4$\\
\textbf{Nós B,C,D,E:} $k = 1$

\vspace{0.2cm}

\importante{
Grau médio: $\displaystyle\langle k \rangle = \frac{1}{n} \sum_{i=1}^{n} k_i = \frac{2|E|}{|V|}$
}
\end{columns}
\end{frame}

\begin{frame}[shrink=10]{Caminhos e Ciclos}
\begin{columns}
\column{0.5\textwidth}
\definicao{Caminho}{
Sequência de vértices $v_1, v_2, \ldots, v_k$ onde existe aresta entre $v_i$ e $v_{i+1}$.
\begin{itemize}
    \item \textbf{Comprimento:} número de arestas
    \item \textbf{Caminho simples:} sem repetição de vértices
\end{itemize}
}

\definicao{Ciclo}{
Caminho fechado: $v_1 = v_k$
}

\column{0.5\textwidth}
\begin{center}
\begin{tikzpicture}[scale=0.9]
    \node[circle, draw, fill=azulclaro!30] (1) at (0,0) {1};
    \node[circle, draw, fill=azulclaro!30] (2) at (2,0) {2};
    \node[circle, draw, fill=azulclaro!30] (3) at (2,2) {3};
    \node[circle, draw, fill=azulclaro!30] (4) at (0,2) {4};
    \node[circle, draw, fill=azulclaro!30] (5) at (1,1) {5};

    \draw[thick, red, line width=1.5pt] (1) -- (2);
    \draw[thick, red, line width=1.5pt] (2) -- (3);
    \draw[thick, red, line width=1.5pt] (3) -- (4);

    \draw[thick] (4) -- (1);
    \draw[thick] (1) -- (5);
    \draw[thick] (2) -- (5);
    \draw[thick] (3) -- (5);
    \draw[thick] (4) -- (5);

    \node[red] at (1, -0.8) {Caminho: 1-2-3-4};
\end{tikzpicture}
\end{center}
\end{columns}

\vspace{0.3cm}

\begin{exampleblock}{Importância Biológica}
\begin{itemize}
    \item Caminhos: vias de sinalização, rotas metabólicas
    \item Ciclos: feedback loops, ciclos metabólicos
\end{itemize}
\end{exampleblock}
\end{frame}

\begin{frame}[shrink=20]{Componentes Conectados}
\definicao{Componente Conectado}{
Subconjunto máximo de vértices tal que existe um caminho entre qualquer par de vértices.
}

\vspace{0.3cm}

\begin{center}
\begin{tikzpicture}[scale=0.9]
    % Componente 1
    \node[circle, draw, fill=azulclaro!30] (A) at (0,0) {A};
    \node[circle, draw, fill=azulclaro!30] (B) at (1,1) {B};
    \node[circle, draw, fill=azulclaro!30] (C) at (1,-1) {C};
    \draw[thick] (A) -- (B);
    \draw[thick] (B) -- (C);
    \draw[thick] (C) -- (A);

    % Componente 2
    \node[circle, draw, fill=verdeacido!30] (D) at (3,0) {D};
    \node[circle, draw, fill=verdeacido!30] (E) at (4,0) {E};
    \draw[thick] (D) -- (E);

    % Componente 3 (isolado)
    \node[circle, draw, fill=laranjacel!30] (F) at (6,0) {F};

    \node[below] at (0.5,-2) {Componente 1};
    \node[below] at (3.5,-2) {Componente 2};
    \node[below] at (6,-2) {Componente 3};
\end{tikzpicture}
\end{center}

\importante{
O maior componente conectado é chamado de \textbf{componente gigante}. Em redes biológicas, geralmente contém a maioria dos nós.
}
\end{frame}

\begin{frame}[shrink=15]{Diâmetro e Comprimento Médio de Caminho}
\definicao{Distância entre Vértices}{
$d(u, v)$ = comprimento do menor caminho entre $u$ e $v$ (geodésica)
}

\vspace{0.3cm}

\begin{columns}
\column{0.5\textwidth}
\definicao{Diâmetro}{
Maior distância entre qualquer par de vértices:
\begin{equation*}
D = \max_{u,v \in V} d(u, v)
\end{equation*}
}

\definicao{Caminho Médio}{
\begin{equation*}
L = \frac{1}{n(n-1)} \sum_{u \neq v} d(u, v)
\end{equation*}
}

\column{0.5\textwidth}
\begin{exampleblock}{Propriedade Small-World}
Muitas redes biológicas têm:
\begin{itemize}
    \item Diâmetro pequeno
    \item $L \propto \log(n)$
    \item Alcance rápido entre nós
\end{itemize}
\end{exampleblock}

\begin{alertblock}{Fenômeno Six Degrees}
Na maioria das redes sociais e biológicas, $L \approx 6$
\end{alertblock}
\end{columns}
\end{frame}

\begin{frame}[fragile, shrink=5]{Criando e Analisando Grafos com NetworkX (1/2)}
\begin{lstlisting}[language=Python, caption=Exemplo básico com NetworkX, basicstyle=\ttfamily\footnotesize]
import networkx as nx
import matplotlib.pyplot as plt

# Criar grafo nao-direcionado
G = nx.Graph()

# Adicionar nos e arestas
G.add_nodes_from([1, 2, 3, 4, 5])
G.add_edges_from([(1,2), (1,3), (2,3), (3,4), (4,5)])

# Propriedades basicas
print(f"Numero de nos: {G.number_of_nodes()}")
print(f"Numero de arestas: {G.number_of_edges()}")
\end{lstlisting}
\end{frame}

\begin{frame}[fragile, shrink=5]{Criando e Analisando Grafos com NetworkX (2/2)}
\begin{lstlisting}[language=Python, basicstyle=\ttfamily\footnotesize]
# Calcular distancias
print(f"Diametro: {nx.diameter(G)}")
print(f"Caminho medio: {nx.average_shortest_path_length(G)}")

# Visualizar
nx.draw(G, with_labels=True, node_color='lightblue',
        node_size=500, font_size=16)
plt.show()
\end{lstlisting}
\end{frame}

\begin{frame}[fragile, shrink=10]{Análise de Matriz de Adjacência (1/2)}
\begin{lstlisting}[language=Python, basicstyle=\ttfamily\footnotesize]
import numpy as np
import networkx as nx

# Criar grafo de exemplo
G = nx.Graph()
G.add_edges_from([(1,2), (1,3), (2,3), (3,4)])

# Obter matriz de adjacencia
A = nx.adjacency_matrix(G).todense()
print("Matriz de Adjacencia:")
print(A)

# Calcular graus a partir da matriz
degrees = np.sum(A, axis=1)
print(f"\nGraus dos vertices: {degrees.T}")
\end{lstlisting}
\end{frame}

\begin{frame}[fragile, shrink=10]{Análise de Matriz de Adjacência (2/2)}
\begin{lstlisting}[language=Python, basicstyle=\ttfamily\footnotesize]
# Matriz de adjacencia elevada ao quadrado
A2 = np.linalg.matrix_power(A, 2)
print("\nMatriz A^2 (caminhos de comprimento 2):")
print(A2)

# Centralidade de autovetor
eigenvalues, eigenvectors = np.linalg.eig(A)
principal_eigenvector = eigenvectors[:, 0]
print(f"\nCentralidade de autovetor: {principal_eigenvector}")
\end{lstlisting}
\end{frame}

% ============================================================================
% SEÇÃO 3: PROPRIEDADES TOPOLÓGICAS DE REDES
% ============================================================================
\section{Propriedades Topológicas de Redes}

\begin{frame}[shrink=10]{Distribuição de Graus}
\definicao{Distribuição de Graus}{
$P(k)$ = fração de nós na rede com grau $k$
}

\vspace{0.3cm}

\begin{columns}
\column{0.5\textwidth}
\textbf{Tipos de Distribuição:}

\begin{exampleblock}{Rede Aleatória (Erdős-Rényi)}
\begin{itemize}
    \item Distribuição de Poisson
    \item $P(k) = \frac{\langle k \rangle^k e^{-\langle k \rangle}}{k!}$
    \item A maioria dos nós tem grau próximo a $\langle k \rangle$
\end{itemize}
\end{exampleblock}

\column{0.5\textwidth}
\begin{exampleblock}{Rede Scale-Free}
\begin{itemize}
    \item Lei de potência (power-law)
    \item $P(k) \sim k^{-\gamma}$
    \item $\gamma$ tipicamente entre 2 e 3
    \item Presença de hubs (nós altamente conectados)
\end{itemize}
\end{exampleblock}
\end{columns}

\vspace{0.3cm}

\importante{
A maioria das redes biológicas segue distribuição scale-free!
}
\end{frame}

\begin{frame}[shrink=15]{Visualizando Distribuições de Grau}
\begin{columns}
\column{0.5\textwidth}
\textbf{Rede Aleatória}
\begin{center}
\begin{tikzpicture}[scale=0.5]
    \draw[->] (0,0) -- (4,0) node[right] {$k$};
    \draw[->] (0,0) -- (0,3) node[above] {$P(k)$};
    \draw[thick, azulescuro, smooth] plot coordinates {
        (0.5,0.1) (1,1) (1.5,2.5) (2,2.8) (2.5,2) (3,0.8) (3.5,0.2)
    };
    \draw[dashed] (2,0) -- (2,2.8);
\end{tikzpicture}
\end{center}

\column{0.5\textwidth}
\textbf{Rede Scale-Free}
\begin{center}
\begin{tikzpicture}[scale=0.5]
    \draw[->] (0,0) -- (4,0) node[right] {$k$};
    \draw[->] (0,0) -- (0,3) node[above] {$P(k)$};
    \draw[thick, verdeacido, domain=1.1:3.8, smooth] plot (\x, {3/\x});
\end{tikzpicture}
\end{center}
\end{columns}

\begin{exampleblock}{Identificando Power-Law}
Gráfico log-log: $\log P(k) = -\gamma \log k + c$ (linha reta)
\end{exampleblock}
\end{frame}

\begin{frame}[shrink=10]{Coeficiente de Agrupamento (Clustering)}
\definicao{Clustering Local}{
$C_i = 2E_i / k_i(k_i - 1)$
}

\begin{columns}
\column{0.5\textwidth}
\begin{center}
\begin{tikzpicture}[scale=0.5]
    \node[circle, draw, fill=laranjacel!50] (A) at (1,1.5) {A};
    \node[circle, draw, fill=azulclaro!30] (B) at (0,0) {B};
    \node[circle, draw, fill=azulclaro!30] (C) at (2,0) {C};
    \node[circle, draw, fill=azulclaro!30] (D) at (1,3) {D};
    \draw[thick] (A) -- (B); \draw[thick] (A) -- (C); \draw[thick] (A) -- (D);
    \draw[thick, verdeacido] (B) -- (C);
    \node[below] at (1,-0.5) {$C_A = 1/3$};
\end{tikzpicture}
\end{center}

\column{0.5\textwidth}
\definicao{Global}{
$C = \frac{1}{n} \sum C_i$
}
\importante{Indica modularidade!}
\end{columns}
\end{frame}

\begin{frame}[shrink=15]{Redes Small-World (Watts-Strogatz)}
\begin{columns}
\column{0.5\textwidth}
\textbf{Características:}
\begin{itemize}
    \item Alto coeficiente de agrupamento (como rede regular)
    \item Pequeno caminho médio (como rede aleatória)
    \item $C \gg C_{random}$
    \item $L \approx L_{random}$
\end{itemize}

\vspace{0.3cm}

\begin{exampleblock}{Modelo WS}
1. Começar com anel regular\\
2. Reconectar cada aresta com probabilidade $p$\\
3. $p \in [0,1]$: regular $\to$ aleatória
\end{exampleblock}

\column{0.5\textwidth}
\begin{center}
\begin{tikzpicture}[scale=0.7]
    % Anel regular
    \foreach \i in {1,...,8} {
        \node[circle, draw, fill=azulclaro!30] (n\i) at ({45*\i-45}:1.5) {};
    }
    \foreach \i [evaluate=\i as \j using {int(mod(\i,8)+1)}] in {1,...,8} {
        \draw[thick] (n\i) -- (n\j);
    }

    % Algumas reconexoes
    \draw[thick, verdeacido, dashed] (n1) -- (n5);
    \draw[thick, verdeacido, dashed] (n3) -- (n7);

    \node[above] at (0,2) {Small-World};
\end{tikzpicture}
\end{center}
\end{columns}

\vspace{0.3cm}

\importante{
Redes small-world permitem comunicação eficiente com conexões locais e alguns atalhos globais.
}
\end{frame}

\begin{frame}[shrink=15]{Redes Scale-Free (Barabási-Albert)}
\begin{block}{Modelo de Ligação Preferencial (Preferential Attachment)}
\begin{enumerate}
    \item Começar com pequeno número de nós $m_0$
    \item Adicionar novo nó com $m$ arestas
    \item Conectar a nós existentes com probabilidade:
    \begin{equation*}
    \Pi(k_i) = \frac{k_i}{\sum_j k_j}
    \end{equation*}
    \item ``Rico fica mais rico'' (hubs atraem mais conexões)
\end{enumerate}
\end{block}

\vspace{0.3cm}

\begin{columns}
\column{0.5\textwidth}
\textbf{Resultados:}
\begin{itemize}
    \item $P(k) \sim k^{-3}$
    \item Hubs naturalmente emergem
    \item Rede sem escala característica
\end{itemize}

\column{0.5\textwidth}
\begin{exampleblock}{Exemplos Biológicos}
\begin{itemize}
    \item Redes de interação proteica
    \item Redes metabólicas
    \item Redes de regulação gênica
    \item Internet, redes sociais
\end{itemize}
\end{exampleblock}
\end{columns}
\end{frame}

\begin{frame}[shrink=20]{Motivos de Rede (Network Motifs)}
\definicao{Motivo de Rede}{
Padrões de interconexão recorrentes que aparecem significativamente mais frequentes que em redes aleatórias.
}

\vspace{0.3cm}

\begin{center}
\begin{tikzpicture}[scale=0.9]
    % Feed-forward loop
    \node (t1) at (0,0) {Feed-forward Loop};
    \node[circle, draw, fill=azulclaro!30] (A1) at (0,-0.8) {X};
    \node[circle, draw, fill=azulclaro!30] (B1) at (-1,-2.2) {Y};
    \node[circle, draw, fill=azulclaro!30] (C1) at (1,-2.2) {Z};
    \draw[->, thick] (A1) -- (B1);
    \draw[->, thick] (A1) -- (C1);
    \draw[->, thick] (B1) -- (C1);

    % Bi-fan
    \node (t2) at (4.5,0) {Bi-fan};
    \node[circle, draw, fill=verdeacido!30] (A2) at (3.8,-0.8) {X};
    \node[circle, draw, fill=verdeacido!30] (B2) at (5.2,-0.8) {Y};
    \node[circle, draw, fill=verdeacido!30] (C2) at (3.8,-2.2) {Z};
    \node[circle, draw, fill=verdeacido!30] (D2) at (5.2,-2.2) {W};
    \draw[->, thick] (A2) -- (C2);
    \draw[->, thick] (A2) -- (D2);
    \draw[->, thick] (B2) -- (C2);
    \draw[->, thick] (B2) -- (D2);

    % Feedback loop
    \node (t3) at (8.5,0) {Feedback Loop};
    \node[circle, draw, fill=laranjacel!30] (A3) at (7.8,-1.5) {X};
    \node[circle, draw, fill=laranjacel!30] (B3) at (9.2,-1.5) {Y};
    \draw[->, thick, bend left=30] (A3) to (B3);
    \draw[->, thick, bend left=30] (B3) to (A3);
\end{tikzpicture}
\end{center}

\begin{exampleblock}{Funções Biológicas}
\begin{itemize}
    \item \textbf{Feed-forward:} filtro de ruído, detecção de persistência
    \item \textbf{Feedback positivo:} switch biestável, memória
    \item \textbf{Feedback negativo:} homeostase, oscilações
\end{itemize}
\end{exampleblock}
\end{frame}

\begin{frame}[shrink=20]{Nós Hub e Vulnerabilidade}
\begin{columns}
\column{0.5\textwidth}
\definicao{Hub}{
Nó com grau muito maior que a média da rede ($k_i \gg \langle k \rangle$)
}

\vspace{0.3cm}

\textbf{Propriedades:}
\begin{itemize}
    \item Centrais para conectividade
    \item Alvos de drogas e terapias
    \item Remoção causa fragmentação
    \item Evolutivamente conservados
\end{itemize}

\column{0.5\textwidth}
\begin{center}
\begin{tikzpicture}[scale=0.8]
    % Hub central
    \node[circle, draw, fill=red!50, minimum size=1cm] (H) at (0,0) {Hub};

    % Nos perifericos
    \foreach \i in {1,...,8} {
        \node[circle, draw, fill=azulclaro!30, inner sep=2pt] (n\i) at ({45*\i-22.5}:2) {};
        \draw[thick] (H) -- (n\i);
    }

    % Algumas conexoes entre perifericos
    \draw[thick] (n1) -- (n2);
    \draw[thick] (n4) -- (n5);
    \draw[thick] (n7) -- (n8);
\end{tikzpicture}
\end{center}
\end{columns}

\vspace{0.3cm}

\importante{
Redes scale-free são robustas a falhas aleatórias mas vulneráveis a ataques direcionados em hubs!
}
\end{frame}

\begin{frame}[fragile, shrink=30]{Calculando Propriedades de Rede com NetworkX}
\begin{lstlisting}[language=Python, basicstyle=\ttfamily\footnotesize]
import networkx as nx
import numpy as np

# Criar rede scale-free (Barabasi-Albert)
G = nx.barabasi_albert_graph(n=100, m=2)

# Analise
degrees = [G.degree(n) for n in G.nodes()]
print(f"Grau medio: {np.mean(degrees):.2f}")

# Clustering e Caminho Medio
C_global = nx.average_clustering(G)
print(f"Clustering global: {C_global:.3f}")

if nx.is_connected(G):
    L = nx.average_shortest_path_length(G)
    print(f"Caminho medio: {L:.3f}")

# Identificar hubs (nos com grau > 2*media)
threshold = 2 * np.mean(degrees)
hubs = [n for n in G.nodes() if G.degree(n) > threshold]
print(f"Hubs encontrados: {len(hubs)}")
\end{lstlisting}
\end{frame}

\begin{frame}[fragile, shrink=30]{Gerando Diferentes Modelos de Rede (1/2)}
\begin{lstlisting}[language=Python, basicstyle=\ttfamily\footnotesize]
import networkx as nx

# 1. Rede Aleatoria (Erdos-Renyi)
G_random = nx.erdos_renyi_graph(n=50, p=0.1)

# 2. Rede Small-World (Watts-Strogatz)
G_sw = nx.watts_strogatz_graph(n=50, k=4, p=0.3)

# 3. Rede Scale-Free (Barabasi-Albert)
G_sf = nx.barabasi_albert_graph(n=50, m=2)
\end{lstlisting}
\end{frame}

\begin{frame}[fragile, shrink=10]{Gerando Diferentes Modelos de Rede (2/2)}
\begin{lstlisting}[language=Python, basicstyle=\ttfamily\footnotesize]
# Comparar propriedades
models = [G_random, G_sw, G_sf]
names = ['Random', 'Small-World', 'Scale-Free']

for G, name in zip(models, names):
    C = nx.average_clustering(G)
    L = nx.average_shortest_path_length(G) if nx.is_connected(G) else float('inf')
    degrees = [d for n, d in G.degree()]

    print(f"\n{name}: Clustering: {C:.3f}, Path: {L:.3f}")
\end{lstlisting}
\end{frame}

% ============================================================================
% SEÇÃO 4: REDES BIOLÓGICAS
% ============================================================================
\section{Redes Biológicas}

\begin{frame}[shrink=5]{Redes de Interação Proteína-Proteína (1/2)}
\begin{columns}
\column{0.5\textwidth}
\textbf{Características:}
\begin{itemize}
    \item Nós: proteínas
    \item Arestas: interações físicas
    \item Não-direcionadas
    \item Scale-free
    \item Modularidade alta
\end{itemize}

\vspace{0.2cm}

\textbf{Detecção Experimental:}
\begin{itemize}
    \item Yeast two-hybrid (Y2H)
    \item Affinity purification MS
    \item Co-immunoprecipitation
\end{itemize}

\column{0.5\textwidth}
\begin{center}
\begin{tikzpicture}[scale=0.7]
    % Proteinas principais
    \node[circle, draw, fill=azulclaro!50, font=\small] (P1) at (0,0) {P1};
    \node[circle, draw, fill=azulclaro!50, font=\small] (P2) at (2,1) {P2};
    \node[circle, draw, fill=azulclaro!50, font=\small] (P3) at (2,-1) {P3};
    \node[circle, draw, fill=azulclaro!50, font=\small] (P4) at (4,0) {P4};
    \node[circle, draw, fill=verdeacido!50, font=\small] (P5) at (3,2) {P5};
    \node[circle, draw, fill=verdeacido!50, font=\small] (P6) at (5,2) {P6};

    % Interacoes
    \draw[thick] (P1) -- (P2);
    \draw[thick] (P1) -- (P3);
    \draw[thick] (P2) -- (P4);
    \draw[thick] (P3) -- (P4);
    \draw[thick] (P2) -- (P5);
    \draw[thick] (P5) -- (P6);
    \draw[thick] (P4) -- (P6);

    \node[below] at (2,-2) {Rede PPI};
\end{tikzpicture}
\end{center}
\end{columns}
\end{frame}

\begin{frame}[shrink=5]{Redes de Interação Proteína-Proteína (2/2)}
\begin{block}{Aplicações Principais}
\begin{itemize}
    \item Predição de função proteica por homologia de interação
    \item Identificação de complexos proteicos
    \item Descoberta de alvos terapêuticos
\end{itemize}
\end{block}

\vspace{0.3cm}

\begin{exampleblock}{Bases de Dados Principais}
\begin{itemize}
    \item \textbf{STRING:} string-db.org — score de confiança integrado
    \item \textbf{BioGRID:} thebiogrid.org — interações experimentais
    \item \textbf{IntAct:} ebi.ac.uk/intact — curada manualmente
\end{itemize}
\end{exampleblock}
\end{frame}

\begin{frame}[shrink=10]{Redes Metabólicas}
\begin{block}{Duas Representações Principais}
\begin{enumerate}
    \item \textbf{Rede de metabólitos:} nós = metabólitos, arestas = reações
    \item \textbf{Rede de reações:} nós = reações, arestas = metabólitos compartilhados
\end{enumerate}
\end{block}

\begin{center}
\begin{tikzpicture}[scale=1.3,
    metabolite/.style={circle, draw, fill=verdeacido!30, minimum size=0.8cm},
    reaction/.style={rectangle, draw, fill=laranjacel!30, minimum size=0.6cm}]

    % Metabolitos
    \node[metabolite] (A) at (0,0) {A};
    \node[metabolite] (B) at (2,0) {B};
    \node[metabolite] (C) at (4,0) {C};
    \node[metabolite] (D) at (2,-2) {D};

    % Reacoes (enzimas)
    \node[reaction] (E1) at (1,0) {E1};
    \node[reaction] (E2) at (3,0) {E2};
    \node[reaction] (E3) at (2,-1) {E3};

    % Fluxos
    \draw[->, thick] (A) -- (E1);
    \draw[->, thick] (E1) -- (B);
    \draw[->, thick] (B) -- (E2);
    \draw[->, thick] (E2) -- (C);
    \draw[->, thick] (B) -- (E3);
    \draw[->, thick] (E3) -- (D);
\end{tikzpicture}
\end{center}
\end{frame}

\begin{frame}[shrink=10]{Redes Regulatórias Gênicas (GRN)}
\begin{columns}
\column{0.5\textwidth}
\textbf{Componentes:}
\begin{itemize}
    \item Genes alvos, TFs, Elementos (promotores)
\end{itemize}

\textbf{Propriedades:}
\begin{itemize}
    \item Direcionadas, Hierárquicas
    \item Motivos regulatórios, FFLs
\end{itemize}

\column{0.5\textwidth}
\begin{center}
\begin{tikzpicture}[scale=0.6,
    gene/.style={rectangle, draw, fill=azulclaro!30, rounded corners},
    tf/.style={ellipse, draw, fill=roxodna!30}]
    \node[tf] (TF1) at (0,2) {TF1};
    \node[tf] (TF2) at (2,2) {TF2};
    \node[gene] (G1) at (0,0) {G-A};
    \node[gene] (G2) at (2,0) {G-B};
    \draw[->, thick, verdeacido] (TF1) -- (G1) node[midway, left] {+};
    \draw[->, thick, red] (TF1) -- (G2) node[midway, left] {-};
    \draw[->, thick, verdeacido] (TF2) -- (G2) node[midway, right] {+};
\end{tikzpicture}
\end{center}
\end{columns}
\end{frame}

\begin{frame}[shrink=10]{Redes de Sinalização Celular (1/2)}
\begin{block}{Características}
\begin{itemize}
    \item Transmissão de sinais extracelulares para resposta intracelular
    \item Cascatas de fosforilação (quinases) e amplificação
    \item Cross-talk entre vias, direcionadas e dinâmicas
\end{itemize}
\end{block}
\end{frame}

\begin{frame}[shrink=10]{Redes de Sinalização Celular (2/2)}
\begin{center}
\begin{tikzpicture}[scale=0.6, node distance=1cm,
    protein/.style={rectangle, draw, fill=azulclaro!30, rounded corners, minimum width=1.5cm}]
    \node[protein] (R) at (0,3) {Receptor};
    \node[protein, below=of R] (K1) {Quinase 1};
    \node[protein, below=of K1] (K2) {Quinase 2};
    \node[protein, below=of K2] (TF) {TF};
    \node[left=0.5cm of R, fill=laranjacel!50, circle] (L) {Ligante};
    \draw[->, thick] (L) -- (R);
    \draw[->, thick] (R) -- (K1);
    \draw[->, thick] (K1) -- (K2) node[midway, right, font=\tiny] {P};
    \draw[->, thick] (K2) -- (TF) node[midway, right, font=\tiny] {P};
    \node[below=0.5cm of TF] {$\Downarrow$ Núcleo};
\end{tikzpicture}
\end{center}
\end{frame}

\begin{frame}[shrink=5]{Bases de Dados: PPI e Metabolismo}
\begin{columns}
\column{0.5\textwidth}
\begin{block}{Redes PPI}
\begin{itemize}
    \item \textbf{STRING:} string-db.org
    \item \textbf{BioGRID:} thebiogrid.org
    \item \textbf{IntAct:} ebi.ac.uk/intact
\end{itemize}
\end{block}

\column{0.5\textwidth}
\begin{block}{Redes Metabólicas}
\begin{itemize}
    \item \textbf{KEGG:} genome.jp/kegg
    \item \textbf{Reactome:} reactome.org
    \item \textbf{BioCyc:} biocyc.org
\end{itemize}
\end{block}
\end{columns}
\end{frame}

\begin{frame}[shrink=10]{Bases de Dados: Regulação e Sinalização}
\begin{columns}
\column{0.5\textwidth}
\begin{block}{Redes Regulatórias}
\begin{itemize}
    \item \textbf{RegulonDB:} E. coli
    \item \textbf{TRANSFAC:} TFs
    \item \textbf{ENCODE:} Elementos
\end{itemize}
\end{block}

\column{0.5\textwidth}
\begin{block}{Redes de Sinalização}
\begin{itemize}
    \item \textbf{KEGG Pathway}
    \item \textbf{NetPath:} vias de sinalização
    \item \textbf{SignaLink:} cross-talk
\end{itemize}
\end{block}
\end{columns}

\vspace{0.3cm}

\importante{
Integração de múltiplas bases de dados permite análise multi-ômica!
}
\end{frame}

\begin{frame}[fragile, shrink=10]{Acessando Bases de Dados com Python (1/2)}
\begin{lstlisting}[language=Python, basicstyle=\ttfamily\footnotesize]
import requests
import networkx as nx

# Exemplo: Rede PPI do STRING
species, gene, thresh = '9606', 'TP53', 700
url = f'https://string-db.org/api/json/network?identifiers={gene}&species={species}&required_score={thresh}'
response = requests.get(url)
data = response.json()
\end{lstlisting}
\end{frame}

\begin{frame}[fragile, shrink=10]{Acessando Bases de Dados com Python (2/2)}
\begin{lstlisting}[language=Python, basicstyle=\ttfamily\footnotesize]
# Construir grafo
G = nx.Graph()
for interaction in data:
    n1, n2 = interaction['preferredName_A'], interaction['preferredName_B']
    G.add_edge(n1, n2, weight=interaction['score'])

print(f"Rede {gene}: {G.number_of_nodes()} nos, {G.number_of_edges()} arestas")
\end{lstlisting}
\end{frame}

\begin{frame}[shrink=20]{Exemplo: Rede de Interação de TP53}
\begin{columns}
\column{0.5\textwidth}
\textbf{Gene TP53 (p53):}
\begin{itemize}
    \item Supressor tumoral
    \item ``Guardião do genoma''
    \item Hub na rede PPI humana
    \item Interage com $>100$ proteínas
    \item Mutado em $>50\%$ dos cânceres
\end{itemize}

\vspace{0.3cm}

\textbf{Principais Parceiros:}
\begin{itemize}
    \item MDM2 (regulação negativa)
    \item ATM, ATR (sensores de dano)
    \item p21 (ciclo celular)
    \item BAX, PUMA (apoptose)
\end{itemize}

\column{0.5\textwidth}
\begin{center}
\begin{tikzpicture}[scale=0.7]
    % TP53 como hub central
    \node[circle, draw, fill=red!50, minimum size=1.2cm, font=\small\bfseries] (TP53) at (0,0) {TP53};

    % Parceiros
    \node[circle, draw, fill=azulclaro!30, font=\tiny] (MDM2) at (-2,1.5) {MDM2};
    \node[circle, draw, fill=azulclaro!30, font=\tiny] (ATM) at (-2,-1.5) {ATM};
    \node[circle, draw, fill=azulclaro!30, font=\tiny] (p21) at (2,1.5) {p21};
    \node[circle, draw, fill=azulclaro!30, font=\tiny] (BAX) at (2,-1.5) {BAX};
    \node[circle, draw, fill=azulclaro!30, font=\tiny] (PUMA) at (0,2.2) {PUMA};
    \node[circle, draw, fill=azulclaro!30, font=\tiny] (CDK) at (0,-2.2) {CDK2};

    % Interacoes
    \draw[thick] (TP53) -- (MDM2);
    \draw[thick] (TP53) -- (ATM);
    \draw[thick] (TP53) -- (p21);
    \draw[thick] (TP53) -- (BAX);
    \draw[thick] (TP53) -- (PUMA);
    \draw[thick] (TP53) -- (CDK);

    % Algumas interacoes secundarias
    \draw[thick, gray, dashed] (MDM2) -- (p21);
    \draw[thick, gray, dashed] (p21) -- (CDK);
\end{tikzpicture}
\end{center}
\end{columns}

\importante{
Perda de TP53 desconecta múltiplas vias críticas!
}
\end{frame}

% ============================================================================
% SEÇÃO 5: ANÁLISE DE CONTROLE METABÓLICO
% ============================================================================
\section{Análise de Controle Metabólico}

\begin{frame}[shrink=10]{Introdução à Análise de Controle Metabólico (MCA)}
\definicao{MCA (Metabolic Control Analysis)}{
Teoria quantitativa que analisa como mudanças em parâmetros individuais (atividade enzimática, concentrações) afetam propriedades sistêmicas (fluxos, concentrações).
}

\vspace{0.3cm}

\begin{block}{Questões Fundamentais}
\begin{itemize}
    \item Qual enzima controla o fluxo em uma via metabólica?
    \item Como perturbações locais afetam o sistema global?
    \item Onde devemos intervir terapeuticamente?
    \item Como prever efeitos de mutações ou drogas?
\end{itemize}
\end{block}

\vspace{0.3cm}

\begin{exampleblock}{Desenvolvida por}
Henrik Kacser, Jim Burns (1973) e Reinhart Heinrich, Tom Rapoport (1974)
\end{exampleblock}
\end{frame}

\begin{frame}[shrink=25]{Coeficientes de Controle}
\definicao{Coeficiente de Controle de Fluxo}{
Mede como mudança na atividade da enzima $i$ afeta o fluxo $J$:
\begin{equation*}
C_i^J = \frac{\partial J / J}{\partial E_i / E_i} = \frac{\partial \ln J}{\partial \ln E_i}
\end{equation*}
Elasticidade normalizada do fluxo em relação à enzima.
}

\vspace{0.3cm}

\definicao{Coeficiente de Controle de Concentração}{
Mede como mudança na atividade da enzima $i$ afeta a concentração de metabólito $S$:
\begin{equation*}
C_i^S = \frac{\partial S / S}{\partial E_i / E_i} = \frac{\partial \ln S}{\partial \ln E_i}
\end{equation*}
}

\importante{
Coeficientes de controle são propriedades sistêmicas, não locais!
}
\end{frame}

\begin{frame}[shrink=25]{Coeficientes de Elasticidade}
\definicao{Elasticidade}{
Propriedade local que mede sensibilidade da taxa de reação $v_i$ a mudanças em metabólito $S$:
\begin{equation*}
\varepsilon_S^{v_i} = \frac{\partial v_i / v_i}{\partial S / S} = \frac{\partial \ln v_i}{\partial \ln S}
\end{equation*}
}

\vspace{0.3cm}

\begin{exampleblock}{Interpretação}
\begin{itemize}
    \item $\varepsilon > 0$: $S$ ativa a reação (substrato ou ativador)
    \item $\varepsilon < 0$: $S$ inibe a reação (inibidor)
    \item $|\varepsilon| > 1$: alta sensibilidade (cooperatividade)
    \item $|\varepsilon| < 1$: baixa sensibilidade
\end{itemize}
\end{exampleblock}

\vspace{0.3cm}

\begin{block}{Para Cinética de Michaelis-Menten}
\begin{equation*}
v = \frac{V_{max}S}{K_M + S} \quad \Rightarrow \quad \varepsilon_S^v = \frac{K_M}{K_M + S}
\end{equation*}
\end{block}
\end{frame}

\begin{frame}[shrink=10]{Teoremas de Soma: Concentração}
\begin{block}{Controle de Concentração}
Para metabólito intermediário $S$:
\begin{equation*}
\sum_{i=1}^{n} C_i^S = 0
\end{equation*}
Aumento e diminuição de concentração devem se balancear.
\end{block}

\importante{
Não há uma única enzima limitante! O controle é distribuído.
}
\end{frame}

\begin{frame}[shrink=10]{Teorema da Conectividade}
\begin{block}{Relaciona Controle e Elasticidade}
Para metabólito intermediário $S$:
\begin{equation*}
\sum_{i=1}^{n} C_i^J \varepsilon_S^{v_i} = 0
\end{equation*}

Para reações consecutivas $v_1 \rightarrow S \rightarrow v_2$:
\begin{equation*}
C_1^J \varepsilon_S^{v_1} + C_2^J \varepsilon_S^{v_2} = 0
\end{equation*}
\end{block}

\vspace{0.3cm}

\begin{exampleblock}{Interpretação}
\begin{itemize}
    \item Liga propriedades locais (elasticidades) a globais (controle)
    \item Permite calcular coeficientes de controle a partir de elasticidades
    \item Fundamental para modelagem modular
\end{itemize}
\end{exampleblock}
\end{frame}

\begin{frame}[shrink=15]{Exemplo: Via Linear Simples}
\begin{center}
\begin{tikzpicture}[scale=1.2,
    node distance=2cm,
    metabolite/.style={circle, draw, fill=verdeacido!30, minimum size=0.8cm},
    reaction/.style={->, thick, azulescuro}]

    \node[metabolite] (X0) at (0,0) {$X_0$};
    \node[metabolite] (S1) at (2,0) {$S_1$};
    \node[metabolite] (S2) at (4,0) {$S_2$};
    \node[metabolite] (X1) at (6,0) {$X_1$};

    \draw[reaction] (X0) -- (S1) node[midway, above, font=\small] {$v_1$};
    \draw[reaction] (S1) -- (S2) node[midway, above, font=\small] {$v_2$};
    \draw[reaction] (S2) -- (X1) node[midway, above, font=\small] {$v_3$};

    \node[below=0.5cm of X0] {Fixo};
    \node[below=0.5cm of X1] {Fixo};
\end{tikzpicture}
\end{center}

\vspace{0.3cm}

\textbf{Estado estacionário:} $v_1 = v_2 = v_3 = J$ (fluxo constante)

\vspace{0.3cm}

\begin{block}{Usando Teoremas de MCA}
\begin{align*}
C_1^J + C_2^J + C_3^J &= 1 \quad \text{(Soma)}\\
C_1^J \varepsilon_{S_1}^{v_1} + C_2^J \varepsilon_{S_1}^{v_2} &= 0 \quad \text{(Conectividade para } S_1)\\
C_2^J \varepsilon_{S_2}^{v_2} + C_3^J \varepsilon_{S_2}^{v_3} &= 0 \quad \text{(Conectividade para } S_2)
\end{align*}
\end{block}
\end{frame}

\begin{frame}[shrink=25]{Cálculo de Coeficientes de Controle}
\textbf{Para via linear:} $X_0 \xrightarrow{v_1} S_1 \xrightarrow{v_2} S_2 \xrightarrow{v_3} X_1$

\vspace{0.3cm}

\begin{block}{Sistema de Equações}
Supondo $v_1$ não depende de $S_1$ e $v_3$ não depende de $S_2$:
\begin{align*}
\varepsilon_{S_1}^{v_1} &= 0, \quad \varepsilon_{S_2}^{v_3} = 0\\
C_1^J + C_2^J + C_3^J &= 1\\
C_2^J \varepsilon_{S_1}^{v_2} &= 0 \quad \Rightarrow \quad C_2^J = 0 \text{ ou } \varepsilon_{S_1}^{v_2} = \infty
\end{align*}
\end{block}

\vspace{0.3cm}

\begin{exampleblock}{Se todas elasticidades são similares}
Aproximação: controle distribuído igualmente
\begin{equation*}
C_1^J \approx C_2^J \approx C_3^J \approx \frac{1}{3}
\end{equation*}
Cada enzima contribui $\sim 33\%$ para controle do fluxo.
\end{exampleblock}
\end{frame}

\begin{frame}[shrink=10]{Exemplo Prático: O Início da Glicólise}
\begin{block}{Enunciado do Problema}
Queremos analisar quantitativamente o controle do fluxo na etapa inicial da glicólise em células tumorais (efeito Warburg). O sistema consiste na fosforilação da Glicose ($G$) extracelular pela enzima \textbf{Hexocinase} ($E_1$), produzindo \textbf{Glicose-6-Fosfato} ($S$), que é então isomerizada pela \textbf{G6P Isomerase} ($E_2$) em \textbf{Frutose-6-Fosfato} ($P$).
\end{block}

\vspace{0.1cm}
\begin{center}
$\text{Glicose } (G) \xrightarrow{E_1 \text{ (Hexocinase)}} \text{G6P } (S) \xrightarrow{E_2 \text{ (Isomerase)}} \text{F6P } (P)$
\end{center}
\vspace{0.1cm}

\begin{block}{Parâmetros do Modelo}
\begin{itemize}
    \item Concentração de Glicose mantida fixa por transporte celular: $G = 5.0$ mM.
    \item Hexocinase ($E_1$): $v_1 = \frac{V_{max1} \cdot G}{K_{M1} + G}$, com $V_{max1} = 10.0$ mM/min, $K_{M1} = 1.0$ mM.
    \item Isomerase ($E_2$): $v_2 = \frac{V_{max2} \cdot S}{K_{M2} + S}$, com $V_{max2} = 15.0$ mM/min, $K_{M2} = 2.0$ mM.
\end{itemize}
\end{block}
\end{frame}

\begin{frame}[shrink=10]{Resolução: Passo 1 (Estado Estacionário)}
\begin{block}{Determinar o Fluxo da Via ($J$) e Concentração de G6P ($S$)}
No estado estacionário, as taxas de consumo e produção devem se equilibrar: $v_1 = v_2 = J$.
\end{block}

\begin{enumerate}
    \item \textbf{Cálculo da taxa da Hexocinase ($v_1$)}:
    Como $G$ é fixo, $v_1$ funciona como fluxo de entrada constante:
    \begin{equation*}
    v_1 = \frac{10.0 \cdot 5.0}{1.0 + 5.0} = \frac{50.0}{6.0} \approx 8.33 \text{ mM/min} \quad \Rightarrow \quad J = 8.33 \text{ mM/min}
    \end{equation*}
    
    \item \textbf{Cálculo da concentração estacionária do intermediário G6P ($S$)}:
    \begin{equation*}
    v_2(S) = J \quad \Rightarrow \quad \frac{15.0 \cdot S}{2.0 + S} = \frac{25.0}{3.0}
    \end{equation*}
    Multiplicando em cruz:
    \begin{equation*}
    45.0 \cdot S = 50.0 + 25.0 \cdot S \quad \Rightarrow \quad 20.0 \cdot S = 50.0 \quad \Rightarrow \quad S = 2.50 \text{ mM}
    \end{equation*}
\end{enumerate}

\vspace{0.1cm}
\centering
\textcolor{azulescuro}{\textbf{Conclusão}: O fluxo estacionário da via é $J = 8.33$ mM/min e a concentração de G6P é $S = 2.5$ mM.}
\end{frame}

\begin{frame}{Resolução: Passos 2 e 3 (Elasticidade e Controle)}
\begin{block}{Passo 2: Calcular as Elasticidades Locais}
Mede a sensibilidade de cada reação em relação ao intermediário G6P ($S$):
\begin{itemize}
    \item $\varepsilon_S^{v_1} = 0$ (a Hexocinase não depende de $S$ neste modelo linear simplificado).
    \item $\varepsilon_S^{v_2} = \frac{K_{M2}}{K_{M2} + S} = \frac{2.0}{2.0 + 2.5} = \frac{2.0}{4.5} \approx 0.44$.
\end{itemize}
\end{block}

\begin{block}{Passo 3: Resolver os Teoremas da MCA}
\begin{align*}
C_1^J \varepsilon_S^{v_1} + C_2^J \varepsilon_S^{v_2} &= 0 \quad \text{(Conectividade)} \quad \Rightarrow \quad C_1^J \cdot (0) + C_2^J \cdot (0.44) = 0 \quad \Rightarrow \quad C_2^J = 0 \\
C_1^J + C_2^J &= 1 \quad \text{(Soma)} \quad \Rightarrow \quad C_1^J + 0 = 1 \quad \Rightarrow \quad C_1^J = 1
\end{align*}
\end{block}

\vspace{0.1cm}
\centering
\textcolor{azulescuro}{\textbf{Resultado}: $C_{HK}^J$ (Hexocinase) = $1.0$ e $C_{PGI}^J$ (Isomerase) = $0.0$. A Hexocinase detém todo o controle sobre o fluxo; aumentar a Isomerase não altera o fluxo $J$ da via.}
\end{frame}

\begin{frame}{Feedback alostérico e a Realidade Celular}
\begin{block}{A Realidade Biológica: Inibição por Produto}
In vivo, a Hexocinase é inibida por feedback alostérico pelo seu próprio produto $S$ ($G6P$):
\begin{equation*}
v_1 = \frac{V_{max1} \cdot G}{K_{M1} + G} \cdot \frac{1}{1 + S / K_{I}} \quad \Rightarrow \quad \varepsilon_S^{v_1} < 0 \quad \text{(elasticidade negativa)}
\end{equation*}
\end{block}

\vspace{0.1cm}
\centering\small
\textcolor{azulescuro}{\textbf{Redistribuição do Controle}: A inibição por produto ($\varepsilon_S^{v_1} \neq 0$) redistribui o controle. O Teorema da Conectividade passa a dar:
\begin{equation*}
C_1^J \varepsilon_S^{v_1} + C_2^J \varepsilon_S^{v_2} = 0 \quad \Rightarrow \quad \frac{C_1^J}{C_2^J} = -\frac{\varepsilon_S^{v_2}}{\varepsilon_S^{v_1}} > 0
\end{equation*}
Ambas as enzimas passam a compartilhar o controle de fluxo, ilustrando o princípio central da MCA de que \textbf{o controle metabólico é distribuído}!}
\end{frame}

\begin{frame}[fragile, shrink=10]{Simulando MCA com Python}
\begin{lstlisting}[language=Python, basicstyle=\ttfamily\footnotesize]
import numpy as np
from scipy.optimize import fsolve

# Taxas da Hexocinase (v_hk) e Isomerase (v_pgi)
def v_hk(G, Vmax1=10.0, Km1=1.0): return Vmax1 * G / (Km1 + G)
def v_pgi(S, Vmax2=15.0, Km2=2.0): return Vmax2 * S / (Km2 + S)

# Estado estacionario (v_hk = v_pgi)
G_fixed = 5.0
S_ss = fsolve(lambda S: v_hk(G_fixed) - v_pgi(S), 2.5)[0]
J_ss = v_pgi(S_ss)

print(f"G6P estacionario (S_ss) = {S_ss:.2f} mM")
print(f"Fluxo estacionario (J_ss) = {J_ss:.2f} mM/min")
\end{lstlisting}
\end{frame}

\begin{frame}[fragile, shrink=10]{Cálculo de Coeficientes de Controle (Python)}
\begin{lstlisting}[language=Python, basicstyle=\ttfamily\scriptsize]
def calc_C(idx, delta=0.01):
    def get_J(v_hk_func, v_pgi_func):
        S = fsolve(lambda S: v_hk_func(G_fixed) - v_pgi_func(S), 2.5)[0]
        return v_pgi_func(S)
    
    if idx == 1:
        Jb = get_J(lambda g: v_hk(g, 10.0), v_pgi)
        Jp = get_J(lambda g: v_hk(g, 10.0*(1+delta)), v_pgi)
    else:
        Jb = get_J(v_hk, lambda s: v_pgi(s, 15.0))
        Jp = get_J(v_hk, lambda s: v_pgi(s, 15.0*(1+delta)))
    return ((Jp - Jb) / Jb) / delta

print(f"C_HK^J  = {calc_C(1):.3f}")
print(f"C_PGI^J = {calc_C(2):.3f}")
\end{lstlisting}
\end{frame}

\begin{frame}[shrink=10]{Limitações da MCA Tradicional}
\begin{block}{Desafios da Modelagem Cinética Clássica}
Embora a MCA forneça uma descrição quantitativa rigorosa do controle do metabolismo, sua aplicação prática enfrenta barreiras significativas:
\begin{itemize}
    \item \textbf{Dependência de Parâmetros}: Exige dados cinéticos detalhados (como $K_M$, $V_{max}$, constantes de inibição) para todas as enzimas envolvidas na rede.
    \item \textbf{Dificuldade de Medição}: Muitos desses parâmetros são difíceis ou impossíveis de se medir de forma confiável \textit{in vivo} em condições fisiológicas.
    \item \textbf{Natureza Local}: As conclusões da MCA (elasticidades e coeficientes de controle) são linearizações válidas estritamente em torno de um \textbf{estado estacionário específico}. Perturbações grandes ou transições metabólicas invalidam a análise linear local.
    \item \textbf{Escalabilidade Limitada}: Modelar redes metabólicas em escala genômica (GEMs) com cinética detalhada é computacional e experimentalmente inviável atualmente.
\end{itemize}
\end{block}
\end{frame}

\begin{frame}[shrink=10]{A Abordagem Baseada em Restrições (Palsson)}
\definicao{Abordagem de Palsson (COBRA)}{
Proposta por Bernhard Palsson e colaboradores, a modelagem \textbf{COBRA} (\textit{Constraint-Based Reconstruction and Analysis}) foca na reconstrução de redes metabólicas em escala genômica usando propriedades estequiométricas fundamentais, eliminando a dependência de parâmetros cinéticos individuais.
}

\begin{columns}
\column{0.5\textwidth}
\begin{block}{Princípios das Restrições}
\begin{itemize}
    \item \textbf{Espaço Nulo ($S \cdot v = 0$)}: Define o estado estacionário estrutural baseado estritamente na estequiometria (Matriz $S$).
    \item \textbf{Restrições Físicas}: Delimita o espaço de fluxos permitidos por meio de termodinâmica e capacidade enzimática ($v_{min} \le v \le v_{max}$).
\end{itemize}
\end{block}

\column{0.5\textwidth}
\begin{exampleblock}{Espaço de Soluções}
As restrições físicas definem um \textbf{cone de fluxo poliedral}. O fenótipo metabólico deve operar dentro deste espaço. Ferramentas como o FBA (\textit{Flux Balance Analysis}) buscam soluções ótimas dentro deste cone maximizando objetivos biológicos (ex: produção de biomassa).
\end{exampleblock}
\end{columns}
\end{frame}

\begin{frame}{Comparativo: MCA vs. Abordagem de Palsson}
\begin{table}[h]
\centering
\small
\resizebox{\textwidth}{!}{
\begin{tabular}{p{2.6cm} p{5.2cm} p{5.2cm}}
\toprule
\textbf{Característica} & \textbf{MCA Tradicional} & \textbf{Abordagem de Palsson (COBRA)} \\
\midrule
\textbf{Foco da Análise} & Regulação cinética local e sensibilidade. & Capacidades de fluxo globais e limites estruturais. \\
\textbf{Requisitos de Dados} & Parâmetros cinéticos ($K_M$, $V_{max}$). & Estequiometria (Matriz $S$) e limites de fluxo. \\
\textbf{Escala de Aplicação} & Vias metabólicas pequenas ou isoladas. & Redes em escala genômica (GEMs). \\
\textbf{Fundamento} & EDOs linearizadas localmente. & Álgebra linear (espaço nulo) e prog. linear (FBA). \\
\textbf{Tratamento do Espaço} & Um único ponto de estado estacionário. & Cone de todas as soluções possíveis. \\
\bottomrule
\end{tabular}
}
\end{table}
\end{frame}

\begin{frame}[shrink=10]{Principais Aplicações do FBA}
\definicao{FBA (Flux Balance Analysis)}{
O FBA é a ferramenta central da abordagem baseada em restrições para calcular fluxos metabólicos estacionários sem requerer cinética detalhada.
}

\begin{block}{Previsões e Resultados Possíveis com FBA}
\begin{itemize}
    \item \textbf{Crescimento Celular}: Predição da taxa máxima de crescimento de biomassa sob diferentes fontes de nutrientes.
    \item \textbf{Essencialidade Gênica}: Identificação \textit{in silico} de knockouts de genes ou reações letais (relevante para novos fármacos).
    \item \textbf{Engenharia Metabólica}: Redirecionamento de fluxos para superproduzir subprodutos industriais (ex: etanol, aminoácidos).
    \item \textbf{Trade-offs Fenotípicos}: Mapeamento de relações de troca (envelopes de produção) entre crescimento e produtos.
\end{itemize}
\end{block}

\vspace{0.1cm}
\centering\small
\textcolor{azulescuro}{\textbf{Visão Integrada:} A modelagem COBRA/FBA delimita a fronteira física de capacidades da rede; a MCA descreve como o controle dinâmico opera localmente dentro dela.}
\end{frame}

% ============================================================================
% SEÇÃO 6: EXERCÍCIOS
% ============================================================================
\section{Exercícios}

\begin{frame}[shrink=15]{Exercícios - Parte 1: Teoria dos Grafos}
\begin{block}{Questão 1: Matriz de Adjacência}
Dada a matriz de adjacência abaixo, desenhe o grafo correspondente e calcule:
\begin{equation*}
A = \begin{bmatrix}
0 & 1 & 1 & 0 \\
1 & 0 & 1 & 1 \\
1 & 1 & 0 & 1 \\
0 & 1 & 1 & 0
\end{bmatrix}
\end{equation*}
\begin{enumerate}
    \item O grau de cada vértice, o grau médio e o número de triângulos.
\end{enumerate}
\end{block}

\begin{block}{Questão 2: Caminho e Diâmetro}
Determine a distância entre 1 e 4, o diâmetro e o comprimento médio.
\end{block}
\end{frame}

\begin{frame}[shrink=15]{Exercícios - Parte 2: Propriedades de Rede}
\begin{block}{Questão 3: Coeficiente de Agrupamento}
Para um nó com grau $k=4$ que tem 3 arestas entre vizinhos, calcule $C_i$.
\end{block}

\begin{block}{Questão 4: Distribuição de Graus}
Uma rede de 100 nós tem 60 com grau 2, 30 com grau 3, 8 com grau 5, 2 com grau 10.
Calcule $\langle k \rangle$ e justifique se é aleatória ou scale-free.
\end{block}
\end{frame}

\begin{frame}[shrink=15]{Exercícios - Parte 3: NetworkX}
\begin{block}{Questão 5: Implementação}
Gere uma rede BA com $n=50$, $m=2$. Calcule a distribuição de graus, identifique hubs e compare com rede ER de mesmo tamanho.
\end{block}

\begin{alertblock}{Questão 6: Robustez}
Simule a remoção de nós (aleatória vs. hubs) em rede scale-free. Qual estratégia fragmenta mais rápido?
\end{alertblock}
\end{frame}

\begin{frame}[shrink=15]{Exercícios - Parte 4: MCA}
\begin{block}{Questão 7: Elasticidade}
Para $v = 100S / (10 + S)$, calcule $\varepsilon_S^v$ para $S=5, 10, 50$.
\end{block}

\begin{block}{Questão 8: Teorema da Soma}
Em via com 5 enzimas: $C_1^J = 0.4$, $C_2^J = 0.3$, $C_3^J = 0.2$, $C_4^J = 0.05$. Qual é $C_5^J$?
\end{block}
\end{frame}

\begin{frame}[shrink=20]{Exercícios - Parte 5: Aplicação}
\begin{alertblock}{Projeto Prático}
Escolha uma via metabólica e construa o grafo no NetworkX. Compare previsões topológicas com dados de MCA da literatura.
\end{alertblock}
\end{frame}

% ============================================================================
% SEÇÃO 7: REFERÊNCIAS
% ============================================================================
\section{Referências}

\begin{frame}{Referências Principais}
\begin{thebibliography}{99}
\bibitem{barabasi} Barabási A.L. (2016). \textit{Network Science}. Cambridge University Press. Disponível em: \url{http://networksciencebook.com}

\bibitem{alon} Alon U. (2019). \textit{An Introduction to Systems Biology: Design Principles of Biological Circuits}, 2nd ed. Chapman \& Hall/CRC.

\bibitem{klipp} Klipp E. et al. (2016). \textit{Systems Biology: A Textbook}, 2nd ed. Wiley-VCH.

\bibitem{fell} Fell D. (1997). \textit{Understanding the Control of Metabolism}. Portland Press.

\bibitem{heinrich} Heinrich R., Schuster S. (1996). \textit{The Regulation of Cellular Systems}. Chapman \& Hall.
\end{thebibliography}
\end{frame}

\begin{frame}[shrink=5]{Leitura Complementar}
\begin{block}{Livros}
\begin{itemize}
    \item Newman M.E.J. (2018). \textit{Networks}, 2nd ed. Oxford University Press.
    \item Voit E.O. (2017). \textit{A First Course in Systems Biology}, 2nd ed.
    \item Palsson B.O. (2015). \textit{Systems Biology: Properties of Reconstructed Networks}. Cambridge.
\end{itemize}
\end{block}

\begin{block}{Artigos Fundamentais}
\begin{itemize}
    \item Barabási A.L., Albert R. (1999). Emergence of scaling in random networks. \textit{Science} 286:509-512.
    \item Watts D.J., Strogatz S.H. (1998). Collective dynamics of small-world networks. \textit{Nature} 393:440-442.
    \item Kacser H., Burns J.A. (1973). The control of flux. \textit{Symp. Soc. Exp. Biol.} 27:65-104.
\end{itemize}
\end{block}
\end{frame}

\begin{frame}{Recursos Online}
\begin{block}{Ferramentas e Software}
\begin{itemize}
    \item \textbf{NetworkX:} \url{https://networkx.org}
    \item \textbf{Cytoscape:} \url{https://cytoscape.org} (visualização de redes)
    \item \textbf{igraph:} \url{https://igraph.org} (R e Python)
    \item \textbf{Gephi:} \url{https://gephi.org} (análise interativa)
\end{itemize}
\end{block}

\begin{block}{Bases de Dados}
\begin{itemize}
    \item \textbf{STRING:} \url{https://string-db.org}
    \item \textbf{KEGG:} \url{https://www.genome.jp/kegg}
    \item \textbf{Reactome:} \url{https://reactome.org}
    \item \textbf{BioGRID:} \url{https://thebiogrid.org}
\end{itemize}
\end{block}
\end{frame}

% ============================================================================
% SLIDE FINAL
% ============================================================================
\begin{frame}
\begin{center}
{\Huge Obrigado!}

\vspace{1cm}

{\Large Capítulo 3: Modelos de Redes Estáticas}

\vspace{1cm}

\textbf{Próximo Capítulo:}\\
Modelos de Redes Dinâmicas e Simulação

\vspace{1cm}

\textit{Dúvidas? Perguntas?}
\end{center}
\end{frame}

\end{document}
